\section{Finite blocks and elementary conventions}
\label{sec:preliminaries}

We first reduce Theorem~\ref{thm:main} to uniformly controlled tests on
finite collections of rows. Throughout the proof the shift $\gamma$ is
fixed. We work on $\mathbb T=\mathbb R/\mathbb Z$ with normalized
Lebesgue measure, using periodic lifts to $\mathbb R$ when counting pairs.
In a lifted pair count the first point is counted per unit length and the
second ranges over all its lifts. After the small-width reduction below,
the widths are sufficiently small that the circle and lifted near-pair
counts agree.

\subsection{Reduction to small fixed-mass blocks}

For $v\ge1$ set $r(v)=\psi(v)/v$, and let
\[
 E_v=\bigcup_{z\in\mathbb Z}
 \bigl((z+\gamma)/v-r(v),(z+\gamma)/v+r(v)\bigr)
 \pmod1.
\]
Thus $x\in E_v$ is exactly the approximation inequality in
Theorem~\ref{thm:main}. We discard rows of width zero.

\begin{lemma}\label{pre:large-widths}
If $v r(v)$ does not tend to zero, then $\limsup_v E_v$ has full measure.
\end{lemma}
\begin{proof}
There are a constant $0<c<1/2$ and an unbounded sequence of $v$ for which
$r(v)\ge c/v$. Let $F_v$ be the union of the row intervals of half-width $c/v$.
Its indicator is bounded by one, has mean $2c$, and for every fixed
interval $J\subset\mathbb T$ satisfies
\[
 \int_J\mathbf1_{F_v}=2c|J|+O(1/v).
\]
The error comes only from the two incomplete periods of length $1/v$.
The same convergence holds for a fixed finite union of intervals.
For a measurable set $A$, approximate $A$ in measure by such a union;
the bound $0\le\mathbf1_{F_v}\le1$ gives
$\int_A\mathbf1_{F_v}\to2c|A|$ along this sequence.
Consequently no positive-measure set avoids a whole tail of the $F_v$,
and therefore none avoids a whole tail of the $E_v$. Since the complement
of a limsup is the countable union of these tail-avoiding sets, the
conclusion follows.
\end{proof}

Henceforth suppose $v r(v)\to0$. Replace every positive width by a
dyadic number between half its value and its value, retaining the notation
$r(v)$. This preserves divergence of $\sum_v\phi(v)r(v)$ and can only
shrink the approximation sets. If Theorem~\ref{thm:main} fails, there is a measurable
set $A\subset\mathbb T$ of measure $\mu>0$ that avoids every $E_v$ in
some tail. Indeed the complement of the limsup is the union over $n$ of
$\bigcap_{v\ge n}E_v^c$.

\begin{lemma}\label{pre:blocks}
For every fixed $0<w\le1/2$ there are arbitrarily late, pairwise disjoint
finite sets $\mathcal B$ of positive-width rows such that
\begin{equation}\label{eq:block-mass}
 w\le W:=\sum_{v\in\mathcal B}\phi(v)r(v)\le2w.
\end{equation}
Their least denominator tends to infinity, and all their spatial widths
tend uniformly to zero.
\end{lemma}
\begin{proof}
The individual terms satisfy $\phi(v)r(v)\le v r(v)\to0$.
In a sufficiently late tail each is at most $w$. Starting after the
previous block, stop at the first partial sum reaching $w$; divergence
ensures that this occurs, and the last summand bounds the overshoot by
$w$. The final assertion follows from $r(v)\to0$.
\end{proof}

All subsequent constructions are on one such finite block. Constants may
depend on $\mu$ and on the fixed hierarchy parameters, but not on the
block, its denominators, or its widths. We choose the hierarchy first and
then take $w$ sufficiently small. The last approximation of $A$ by a union
of intervals will occur only after those constants and the fixed bound on
cell loads introduced in Section~\ref{sec:tables} have been chosen.

\subsection{Arithmetic notation and divisor laws}

For positive integers $d,e$, write $(d,e)$ and $[d,e]$ for their gcd and
lcm. Let $\nu_p(n)$ be the exponent of a prime $p$ in $n$, let $\tau(n)$
be the number of positive divisors, and let $\omega(n)$ be the number of
distinct prime divisors. Put
\[
 R(n)=\frac{\phi(n)}n,\qquad R_p=1-\frac1p.
\]
All logarithms are natural unless a dyadic index is explicitly specified.
The symbols $O$, $\ll$, and $\asymp$ have constants uniform in the varying
block data; dependence on fixed parameters will be indicated when needed.

\begin{definition}\label{pre:deficit-law}
A mask on an integer $L\ge1$ is a divisor $d\mid L$. Its deficit law is
\begin{equation}\label{eq:deficit-law}
 \pi_L(d)=\frac{\phi(L/d)}L.
\end{equation}
An exponent in $d$ records the corresponding deficit from $L$ in $L/d$.
\end{definition}

The identity $\sum_{d\mid L}\phi(d)=L$ makes this a probability law.
Euler factorization makes its prime coordinates independent. More
explicitly, if $p^\nu\parallel L$, then
\[
 \Pr(\nu_p(d)=j)=
 \begin{cases}p^{-j}R_p,&0\le j<\nu,\\p^{-\nu},&j=\nu,
 \end{cases}
 \qquad
 \Pr(\nu_p(d)\ge j)=p^{-j}\quad(0\le j\le\nu).
\]
These probabilities weight row denominators. They are not probabilities
that a numerator lies in a residue class. Numerator counts at a fixed mask
will always be justified separately by a complete period or a specified
translation orbit. Other auxiliary probability laws used later are also
explicit finite weights; the shift itself is never randomized.

\begin{lemma}\label{pre:elementary-prime-bounds}
Uniformly for $z\ge2$, $b\ge a\ge2$, and integers $n\ge1$,
\begin{align}
 \sum_{p\le z}\log p&\ll z,
 &\sum_{p\le z}\frac{\log p}{p}&\ll\log(2z),\label{eq:prime-log-sums}\\
 \sum_{a\le p\le b}\frac1p
 &\ll \frac1{\log a}+\log\frac{\log b}{\log a},
 &\prod_{p\le z}R_p^{-1}&\ge\sum_{1\le m\le z}\frac1m.
 \label{eq:prime-reciprocals}
\end{align}
For every fixed $C,\epsilon>0$,
\begin{equation}\label{eq:divisor-small-power}
 \tau(n)^C+R(n)^{-C}\ll_{C,\epsilon}n^\epsilon.
\end{equation}
\end{lemma}
\begin{proof}
The primes in $(N,2N]$ divide $\binom{2N}{N}\le4^N$.
Summing over dyadic intervals gives the first bound in
\eqref{eq:prime-log-sums}; partial summation gives the second. The same
binomial estimate gives $O(t/\log t)$ primes in a dyadic interval at
scale $t$. Decomposing $[a,b]$ into multiplicative intervals and summing
their reciprocal costs gives the first bound in
\eqref{eq:prime-reciprocals}, including the possible endpoint interval.
The Euler product on its right expands into reciprocals of all positive
integers with prime factors at most $z$, including every integer at most
$z$.

For \eqref{eq:divisor-small-power}, at all sufficiently large primes the
factor $(j+1)^C$ is bounded by $p^{\epsilon j/2}$ for every $j\ge1$,
and $R_p^{-C}\le p^{\epsilon/2}$. Each of the finitely many remaining
primes has a bounded supremum after division by its positive exponential
factor $p^{\epsilon j}$, or contributes a fixed constant in the
$R(n)$ bound. Multiplication proves the assertion.
\end{proof}

We will use Cauchy--Schwarz, H\"older, the Chinese remainder theorem,
Markov's inequality, and finite orbit averaging. Size parameters and
vertex totals in independent arithmetic estimates are local to those
estimates; the centre and table notation is introduced only when needed.

\subsection{A tensor kernel for integer families}

The centre construction will use a kernel bound for pairs of integer
families. We derive it from a concentration observation on a single prime
coordinate, retaining the concentration conclusions for the weighted lcm
estimate in Section~\ref{sec:arithmetic}.
The diagonal-concentration argument is an adaptation of Green and
Walker~\cite[Lemma~2.1]{GW2021} and Hauke-Treuer, Vazquez and
Walker~\cite[Lemma~3.2]{HVW2026}. We include the proof, together with
the kernel-total and off-pivot estimates needed below.
Exponent indices in this subsection belong to $\mathbb Z_{\ge0}$.

\begin{lemma}[One-coordinate concentration]\label{ar:pivot}
Fix $1/2<u<1$ and $A\ge1$. For probability vectors $\alpha,\beta$ and
$0<q<1/2$, put
\[
 T_q(\alpha,\beta)=\sum_{i,j}
 q^{|i-j|}A^{\mathbf1_{i\ne j}}(\alpha_i\beta_j)^u.
\]
Then
\[
 T_q(\alpha,\beta)\le 1+O_{u,A}(q^{1/(1-u)}).
\]
Moreover, if a probability array $(m_{ij})$ satisfies
\[
 m_{ij}\le q^{|i-j|}A^{\mathbf1_{i\ne j}}
 (\alpha_i\beta_j)^u,
\]
there is, for sufficiently small $q$, an index $k$ such that
\[
 d:=(1-\alpha_k)+(1-\beta_k)\ll_{u,A}q^{1/(1-u)}.
\]
For this index the mass on $i,j\ne k$ is $O_{u,A}(d^{2u})$, and the mass
on pairs with exactly one index equal to $k$ and the other at distance at
least $a\ge1$ is $O_{u,A}(q^a d^u)$.
\end{lemma}

\begin{proof}
The off-diagonal convolution kernel has $\ell^1$ norm $O_A(q)$, whereas
$\|\alpha^u\|_2,\|\beta^u\|_2\le1$. Hence its contribution is $O_A(q)$.
The diagonal contribution is at most
\[
 \left(\sum_i\alpha_i^{2u}\sum_j\beta_j^{2u}\right)^{1/2}.
\]
If the total is at least one, the latter expression is $1-O_A(q)$.
Since
\[
 \sum_i\alpha_i^{2u}\le(\max_i\alpha_i)^{2u-1},
\]
both vectors have atoms of mass $1-o(1)$. These atoms have the same index:
the actual diagonal sum is at least $1-O_A(q)$, whereas different
dominant indices would bound it by
$O(d_\alpha^u+d_\beta^u+(d_\alpha d_\beta)^u)=o(1)$, with
$d_\alpha,d_\beta$ their off-atom masses.
Call the common index $k$ and define $d$ as in the statement.

The central term is at most $1-c_ud$ when $d$ is small. The single-off
terms are $O_{u,A}(qd^u)$ by geometric summation and Holder's inequality.
More generally, without any support-size assumption,
\[
 \sum_{|j-k|\ge a}q^{|j-k|}\beta_j^u
 \le\left(\sum_{j\ne k}\beta_j\right)^u
 \left(\sum_{|j-k|\ge a}q^{|j-k|/(1-u)}\right)^{1-u}
 \ll_u q^a d^u.
\]
The double-off terms are $O_{u,A}(d^{2u})$ by the convolution bound applied
to the two off-parts. Consequently
\[
 T_q(\alpha,\beta)
 \le 1-c_ud+O_{u,A}(d^{2u}+qd^u).
\]
Absorb $d^{2u}$ using $2u>1$. Maximizing the remaining expression in $d$
proves the asserted excess bound. If a probability array is dominated by
the kernel, its total is one, so the same inequality forces
$d\ll q^{1/(1-u)}$. Restricting the single-off geometric sum to distances
at least $a$ gives its last asserted bound. When the kernel total is less
than one, the first conclusion is immediate.
For $q$ in a compact subinterval $[q_0,1/2)$, the convolution norm is
bounded by $1+2Aq/(1-q)\le1+2A$. Increasing the constant therefore
extends the first bound to all $0<q<1/2$.
\end{proof}

\begin{lemma}[Tensor kernel]\label{ar:tensor}
For $1/2<u<1$, $s>1-u$, and nonnegative sequences $(x_L),(y_M)$ on the
positive integers,
\[
 \sum_{L,M\ge1}
 \frac{(x_Ly_M)^u}{\bigl(L/(L,M)\cdot M/(L,M)\bigr)^s}
 \ll_{u,s}\left(\sum_Lx_L\sum_My_M\right)^u.
\]
\end{lemma}

\begin{proof}
It suffices first to consider finite supports and positive totals, and
normalize both totals to one. At each sufficiently large prime $p$,
Lemma~\ref{ar:pivot} with
$q=p^{-s}$ bounds the optimal constant in the corresponding
one-coordinate estimate by
$1+O_{u,s}(p^{-s/(1-u)})$. Each of the finitely many small primes also has
a bounded optimal constant, directly from the $\ell^1$ convolution bound.

To tensorize, fix the two exponents at the first prime and apply the
remaining-coordinate inequality to the corresponding slices. Its right
side is the product of their slice totals to the power $u$. The
one-coordinate bound then sums those totals. Repeating this argument gives
the product of these one-prime constants. This product converges because
$s/(1-u)>1$. Zero totals are trivial, and monotone convergence removes
the finite-support restriction.
\end{proof}

\subsection{Residue averages for rational grids}

For positive integers $v,w$, consider the two lifted grids
$(z+\xi)/v$ and $(z'+\xi')/w$, with fixed real phases and positive
half-widths $r_1,r_2$. Put
\[
 g=(v,w),\qquad s=w/g,\qquad t=v/g.
\]
The integer determinant $n=sz-tz'$ records their relative positions.
In particular, distance at most $C\max(r_1,r_2)$ restricts $n$ to an
interval of length $O_C(\eta)$, where
$\eta=[v,w]\max(r_1,r_2)$.

\begin{lemma}[Residue averages on determinant cycles]\label{ar:cycle}
Fix divisors $D_1\mid v$, $D_2\mid w$ and a nonnegative function of
$z\bmod D_1,z'\bmod D_2$. At every integer determinant there are $g$
solutions per unit first-point translate; let $A(n)$ be their average
value. Then $A$ has a period dividing $D_1D_2$, and its mean $\overline A$
is the independent uniform residue mean of the function. For any real
interval $J$,
\[
 \sum_{n\in J\cap\mathbb Z}A(n)
 =|J|\overline A+O(D_1D_2\overline A).
\]
If one also requires $n\equiv n_0\pmod Q$ and $(Q,D_1D_2)=1$, the main
term is $|J|\overline A/Q$ with the same error bound.
\end{lemma}

\begin{proof}
Since $(s,t)=1$, all solutions at a fixed determinant differ by $(t,s)$.
Taking the first numerator over one period of length $v$ gives a cycle of
length $g$. A Bezout solution for an increment of $D_1D_2$ in the
determinant preserves both tested residues, proving periodicity. Averaging
also over $n\bmod[v,w]$ enumerates all $vw$ pairs of numerator residues
modulo $v,w$, proving the mean assertion.

A nonnegative periodic sequence of period $P$ and mean $\overline A$ has
total mass $P\overline A$ in one period. Splitting an interval into whole
periods and two end pieces therefore gives an error $O(P\overline A)$.
Take $P=D_1D_2$. Finally, stepping by a number coprime to $P$ permutes its
residues, so the same proof applies on the stated progression.
\end{proof}
